\section{Discussion}
\label{sec:discussion}

\subsection{Mechanism versus mission-level superiority}
The design chain supports a causal account. Ideal receiver knowledge exposes
redundant in-flight retransmission; lastSent/lastAcked dual memory removes that
ambiguity; v2 then reveals over-suppression because one cooldown governs both
innovation and repetition. V3 resolves the semantic conflict by prioritizing
new information and suppressing only refresh. The rate response in EXP10 and
the four signed transitions in EXP11 support mechanistic adaptivity.

That result is not the same as Pareto superiority. A fixed periodic scheduler
cannot adapt its communication effort to within-mission network changes because
its rate is invariant by construction; however, a well-selected fixed rate can
still provide a competitive mission-level accuracy--cost trade-off. P12.5 is
the concrete counterexample: it retains
\MetricElevenPonetwofiveAccuracyPercent{}\% of Causal-v3 accuracy at
\MetricElevenPonetwofiveCostPercent{}\% of its ACK-weighted cost and
\MetricElevenPonetwofiveBroadcastPercent{}\% of its broadcast cost. EXP11
therefore supports the adaptation mechanism, not universal superiority.

\subsection{Why accounting changes the conclusion}
Periodic baselines require no modeled reverse ACK channel, while Causal-v3
does. DATA-only accounting answers whether the policy sends fewer state
packets; $C_w$ answers whether that remains true after pricing feedback; the
broadcast model additionally credits one simultaneous medium transmission
rather than charging each receiver link. K2 and EXP11 show that favorable
DATA or ACK-weighted comparisons can reverse under another legitimate traffic
convention. Claims must therefore name the unit and accounting model.

\subsection{What the references and boundaries mean}
The ideal-feedback and oracle-periodic policies are ideal-feedback,
oracle-information, or efficiency references. They are not bounds on
attainable accuracy or performance. Similarly, the robustness sweeps characterize where
the fixed causal mechanism operates: saturation, controller degree scaling,
unrecoverable blackouts, mass offset, estimator-triggered traffic, and sampled
timestep dependence. They do not establish stability or universal robustness.
